\chapter{系统总体设计}
\label{ch:arch}

四个执行量是闭环的唯一出口。表面温度、腔温和质量从食品侧进入；水位、管温、门、腔压和剩余电流从安全侧进入。两侧不共用一根信号线：品质模型的错误改写不了断电条件，断电条件也回写不了进度后验。家电面板上的“智能”多数把这两侧混写进同一条加热请求，超温切断会随后验一起抖动。品质回路因此只改请求，安全回路只改指令。

产品上这件事对应两块电路。一块读 NTC 或热电偶给出的加热器温度、水箱浮子给出的水位、门微动开关和漏电互感器，越过阈值就切断加热管和蒸汽阀。另一块读表面温度或腔温，用来算“还要不要加热”。两块若共用同一条请求线，进度估计一旦偏高，加热器温度保护也会跟着松；估计一旦偏低，超温切断又可能把已经接近设计线的过程掐断。已经运行的核把请求和指令分成两列。安全投影插在中间，自变量只有硬测量。

\section{闭环结构}

图\ref{fig:loop}给出依存关系。食品与腔体边界产生可安装观测，观测更新决策状态后验；后验同时进入对偶决策和安全投影。信息动作与品质动作在投影之前分开计算，进入同一个执行器接口后写回边界。

\begin{figure}[htbp]
\centering
\includegraphics[width=0.92\textwidth]{architecture.png}
\caption{感知--决策--执行闭环。信息动作与品质动作在投影之前分开，安全层只读硬测量。}
\label{fig:loop}
\end{figure}

图\ref{fig:mloop}把已经运行的核收成五段：边界、十六条控制体、表面观测、探针或跟踪、加热指令。加热指令写回边界。结构方程的路径积分不在这五段里。

\begin{figure}[htbp]
\centering
\includegraphics[width=0.88\textwidth]{M_loop.png}
\caption{已运行闭环的五段。加热指令写回边界；路径积分不在这段顺序里。}
\label{fig:mloop}
\end{figure}

实现中的时间步按这个顺序执行：读表面温度和平均含水，更新热阻代理和冷点估计，形成加热请求，用加热器温度削减请求，再推进十六条控制体。结构方程的路径积分和有限时域优化不在这个顺序里。热阻更新在“是否更接近参照”这一步使用了由真值温差构成的代理，因此方差下降必须连同该参照一起报告。顺序一旦固定，交叉表比较的对象也就固定：四条策略在同一积分器上的停机偏差，执行器接口始终只有四个量。

\section{五层与时间尺度}

五层对应产品上五类已经存在的零件，而不是另造一台机器。

\begin{description}
\item[过程层] 腔体几何、加热、蒸汽、风机、食品几何与初始场。厨房里就是钢板盒子、六根加热管、蒸汽喷嘴、后壁风机和排汽百叶。仿真中由积分器充当。本次积分器没有单独的腔体流动场。
\item[多尺度模型] 把过程压缩为决策状态。压缩必须保留冷点、表面干耗和局部过热。十六条控制体保留冷点（最低温度）和表面温度。面板读数若只给出腔温，这一层就被跳过了。
\item[结构孪生] 保存父节点固定的结构方程和参数后验。本次核没有这一层的独立积分。一次烹饪的短序列上，能量守恒的边保持固定。
\item[对偶决策] 计算信息价值、品质价值，并选择动作。本次对偶策略是固定时间窗加饱和线性请求。窗内蒸汽和风机按时间表离开平衡，窗外加热请求由估计温度、方差和厚度决定。
\item[安全执行] 电气、热、腔压、门封，以及只依赖这些测量的投影。这一层不调用网络。本次只有加热器温度进入削减。水位、门和漏电在规则里同样能把加热写成零。
\end{description}

执行器接口只有一个。动作进入接口之前带有标记：信息、品质或降级保温。核里的模式整数 \(0,1,2\) 就是这个标记的退化形式。保持时请求降到 \(0.15\)，仍然经过同一条加热器更新。种子 7 在 1872 s 切入保持，色带从表面向冷点回落，加热条不再拉满。

四个时间尺度按快慢分层。毫秒级剩余电流切断可以打断分钟级食品估计；秒级探针窗口短于内部热阻时间常数；分钟级品质加热覆盖剩余过程；过程之间的先验写回发生在停机之后。安全层的时钟独立于估计器，慢计算关不掉快保护。家用电热器具上这件事已经成立：漏电互感器不经过食品模型。一体机要保持的，是进度估计同样走不进这条快回路。

\section{决策状态}

操作者需要的判断只有四个：能否停止，表面含水是否偏离目标，温度是否还在散开，继续施加边界是否安全。对应的隐藏量构成决策状态
\begin{equation}
\label{eq:z}
\vect{z}(t)=\big(
T_{\mathrm{cold}},\,
\bar T_s,\,
\sigma_T,\,
\bar\varepsilon_l,\,
\varepsilon_{l,s},\,
\alpha_m,\,
\alpha_{b,s},\,
S_s,\,
m_{\mathrm{loss}}
\big).
\end{equation}
依次为冷点温度、表面均温、温度场标准差、平均液相分数、表面液相分数、致死进度、表面褐变、表面结壳与累计质量损失。停机看的是冷点和致死当量，不看全场的 \(L^2\) 范数。腔温只描述空气。新增一路观测必须降低上列至少一个判断的 \(\Delta\mathcal{I}\)，否则不进入观测向量。

主日志给出 \(T_{\mathrm{cold}}\)、表面温度、平均含水和 \(\alpha_m\)，没有单独给出 \(\sigma_T\) 的决策列。温度离散、表面液相和结壳没有作为结果列输出。九个分量里，停机真正读取的是冷点和致死进度；表面液相和结壳进入路径 B、C 的显示量；质量损失在本次核里几乎不动，与第\ref{ch:ablate}章质量通道开/关零效应同方向。决策状态因此比日志宽、比全场窄：宽的那一侧用来标明还缺哪一路观测，窄的那一侧用来固定交叉表比较的是哪几个量。

操作者在面板上真正需要回答的，仍是那四个判断。冷点有没有越过设计线，决定能不能停；表面含水有没有偏离目标，决定蒸汽该不该关；温度场还在不在散开，决定风机该不该抬；水位、门、漏电和管温有没有越界，决定四个量能不能离开零。九维状态把这四个判断写成同一组符号，食材类别只改先验和终端目标，分量定义保持不变。这是可迁移的工程含义：换一块土豆或一块面团，跟踪器仍读表面温度和平均含水，停机仍看 \(\alpha_m\)。

家用面板很少把这四个判断同时写出来。常见写法是一个温度、一个分钟数、一个档位名称。温度对应腔温或管壁设定，分钟数对应开环菜单，档位名称对应允许集合。冷点进度被折叠进“中等、25 分钟”，表面含水被折叠进“蒸”或“烤”，温度场是否散开被折叠进风机的默认转速，安全四条被折叠进一块看不见的保护板。已经运行的核把这四件事重新拆开。停机读冷点和 \(\alpha_m\)，蒸汽读表面含水是否偏离，风机读温度场是否还在散开，硬切断读水位、门、漏电和管温。面板上的三个字仍然可以保留，它们只限制四个量离开零的许可。

水位开关、门微动、漏电互感器和管壁限温器在拆机时分属四个零件号。品质回路的表面热电偶或红外是第五个零件。五条线若焊到同一块继电器上，进度估计的一次误报会变成一次误断电。已经运行的核把前四条放进\eqref{eq:hard}，第五条只进入跟踪器。家用一体机的电路板上这件事常常已经分开，软件层又把它们重新折进同一条“加热请求”。折回去的那一步，正是超温切断随后验抖动的来源。维修手册里前四个零件各有动作值和复归值；进度估计没有动作值，只有停机门槛随方差和厚度移动。两套数不该写在同一页校准表上。

闭环的五段在产品上也能按零件点名。边界是加热管、喷嘴、风机和百叶；十六条控制体是食品沿厚度的离散；表面观测是门玻璃内侧或搁架上装得上的温度；探针或跟踪是进门后第三分钟那一阵白雾，以及其后的点估计；加热指令写回管壁。操作者拧档位时，改变的是边界段四个量离开零的许可。拧温度旋钮时，若接到腔温比例，改变的是空气温度的闭合，不是冷点进度。拧分钟旋钮时，改变的是开环切断时刻。对偶把分钟旋钮换成停机判据，探针窗占用蒸汽和风机 40 s。五段顺序固定之后，交叉表比较的对象也就固定：四条策略在同一积分器上的停机偏差。

\section{三档只改允许集合}

蒸、烤、组合不增加执行器，只限制四个量的允许集合：
\begin{equation}
\label{eq:mode}
(u,s,f,e)=
\begin{cases}
(\min(u,0.25),\,s,\,f,\,e) & \text{蒸}\\
(u,\,0,\,f,\,\max(e,0.45)) & \text{烤}\\
(u,\,s,\,f,\,e) & \text{组合}.
\end{cases}
\end{equation}
食品模块收到的是通量，不读取档位名称。面板上的三个字对应的是这三条不等式，不是三套独立的传热定律。蒸档把加热请求夹在 0.25 以下，蒸汽仍可离开零，冷凝潜热是这一档的主边界。烤档把蒸汽写成零，并把排汽抬到至少 0.45，表面被推干，冷点只靠导热跟上。组合档放开四个量，蒸汽和干热可以在同一次过程里先后出现。种子 7、肉糜、30 mm、对偶策略下，组合档的停止时刻是 1866 s。烤档的蒸汽列恒为 0。蒸档的加热请求不超过 0.25。

硬切断为
\begin{equation}
\label{eq:hard}
u_{\mathrm{cmd}}=
\begin{cases}
0 & \text{水位、门、漏电或管温触发}\\
u_{\mathrm{req}}\big(1-0.75\,\mathrm{clip}((T_h-230)/40)\big) & T_h>230\\
u_{\mathrm{req}} & \text{其余}.
\end{cases}
\end{equation}
加热器满功率平衡在 \(196.5^\circ\mathrm{C}\)，仿真中 \(230^\circ\mathrm{C}\) 削减不动作。规则仍与后验无关。水位干了仍加热，是干烧；门开着仍喷蒸汽，是烫伤；漏电仍通电，是电击。这三件事不随后验抖动，是因为硬切断的自变量里没有进度估计。

三档允许集合决定的是通量的上界。种子 7、肉糜、30 mm、对偶策略下，组合档在 1866 s 停机，\(e_{\alpha}=+0.145\)，能量 1864；烤档关闭蒸汽后能量升到 2353，进度同样走到上沿；蒸档把加热夹在 0.25，进度停在 0.021，方差留在先验 0.64。同一跟踪器、同一探针窗，改变的是四个量离开零的许可。档位研究里 50 mm 三种档位全部低于 \(-0.2\)：允许集合改不了扩散时间。

闭环每个时间步按这个顺序执行：读表面温度和平均含水，更新热阻代理和冷点估计，形成加热请求，用加热器温度削减请求，再推进十六条控制体。信息动作只占用 \([150,190)\) 这一段，其余时间加热请求由估计温度、方差和厚度决定。安全层插在请求与指令之间。下文的“对偶”指这条固定窗加饱和线性请求，“安全”指只读硬测量的投影。开环把停止时刻写成厚度的函数，腔温比例把停止时刻写成时域终点，无探针和对偶共用饱和线性加热律。交叉表比较的是这四条入口在同一积分器上的停机偏差。

演示页面把同一闭环摊成可见零件，见第\ref{ch:demo}章。六根 U 形加热管对应加热指令，左侧喷嘴与蒸汽粒子对应蒸汽阀，后壁十叶风机对应风机，右上七片百叶对应排汽。水位出现在水箱液面，门出现在铰链开度。九维决策状态里，停机真正读取的冷点和致死进度写在底部曲线的右轴和中栏读数上；表面液相和结壳进入右栏路径 B、C 的显示量。页面读取已经写出的状态流，着色器颜色是两个标量的插值。

四个旋钮在厨房里各有一次被拧动的机会，拧完之后机器按允许集合自己走。加热旋钮抬高的是管壁辐射和对流份额，表面先亮、先干，冷点仍按厚度平方跟上。蒸汽旋钮拉满时，阀侧送出饱和蒸汽，界面上压差可正可负：薄件更容易吃到冷凝，厚件更容易只把表面打湿。排汽旋钮决定湿空气能不能离开，烤档把它的下限写成 0.45，为的是把表面推过结壳门槛；家用组合档把表面湿度维持得较高，S1 里那条门槛常常达不到。风机旋钮改边界层厚度，探针段把它拧到 1.3，\(h\) 从 34 升到 83.5，改的是 \(R_{\mathrm{ext}}\)，不是把肉吹熟。档位旋钮不增加执行器，只把这四个量写成三条不等式。操作者听见的铃声，来自开环时钟、腔温设定、点估计或点估计加短窗。切开之后核对的仍是冷点。

家用电路把这件事做成两块板。电源板上走漏电、管温和水位；主控板上走腔温、分钟和档位。两块板之间若只留一根加热继电器，进度估计的一次误报就会变成一次误断电。已经运行的核把请求写在品质列，把指令写在安全列，投影插在中间。种子 7 的工作点上两列重合，因为加热器停在 \(196.5^\circ\mathrm{C}\)。重合是工作点，不是布线。维修时若把两列焊回一根，故障会重新变成“有时提前切断、有时超温抖动”。侧栏把两列并排写出，演示页上也能看见这一隔离。

面板上的三个字——蒸、烤、组合——对应三条不等式。蒸档把加热请求夹在 0.25 以下，蒸汽可以离开零：冷凝潜热从外向内送，管壁辐射被压住，表面不容易上色。烤档把蒸汽写成零，并把排汽抬到至少 0.45：表面被推干，冷点失去冷凝项。组合档放开四个量，探针窗才有空间把蒸汽拉到 1.0、风机拉到 1.3。种子 7、肉糜、30 mm、对偶策略下，组合档在 1866 s 停机，\(e_{\alpha}=+0.145\)，能量 1864；烤档关闭蒸汽后能量升到 2353，进度同样走到上沿；蒸档进度停在 0.021，方差留在先验 0.64。档位名称写在面板上。允许集合写在四个旋钮上。食品模块收到的始终是通量。

家用面板上的手感把这三条不等式收成两只旋钮加一个档位键。档位键三档停顿，拧到蒸、烤或组合，允许集合当场改写。温度旋钮连续，接到腔温比例时把空气拉住。分钟旋钮按 5 min 一格，接到开环时把切断时刻写成厚度档上的 900、1170 或 1530 s。对偶把分钟旋钮换成停机判据，并在进门后第三分钟占用蒸汽和风机 40 s。操作者听见的铃声，来自这四条入口里被接到面板上的那一条。切开之后核对的仍是冷点。
